\section{Introduction}
\label{sec:introduction}

Vision-language and embodied-language models are increasingly evaluated in
simulated 3D environments where agents must connect text to objects,
viewpoints, and actions. These evaluations often treat the rendered scene as a
stable substrate. In practice, however, assets are often converted across
material systems for scale rendering, distribution, viewer export, or robotics
reuse. A seemingly
engineering-level conversion can therefore become a hidden intervention in the
benchmark itself.

We study this problem in NVIDIA Isaac Sim, where photorealistic assets commonly
use Material Definition Language (MDL) shaders while portable USD and glTF-style
pipelines usually expect UsdPreviewSurface or PBR-like materials. Converting
MDL materials to UsdPreviewSurface can improve interoperability, but it can also
change color, reflectance, lighting response, or texture cues. For
language-grounding benchmarks, the key question is not just whether two renders
look similar, but whether a model still
answers the same object query and points to the same target after the material
representation changes.

This paper treats material conversion as a controlled within-simulation
perturbation. We ask: \emph{when synthetic-scene assets are converted for
toolchain compatibility, what evidence is needed before using the converted
renders or scenes for language-grounding and embodied-data claims?} The answer
is evidence-gated rather than binary: proxy similarity can pass while grounding
coordinates, material mechanisms, and embodied metrics answer different
reliability questions. This matters for embodied-language benchmarks such as
GRScenes~\cite{Wang2024GRUtopia}, where downstream experiments may combine
large USD scenes, object-level goals, VLM prompting, and navigation frameworks
such as InternNav/DualVLN~\cite{Wei2026Ground}. Figure~\ref{fig:acl_method_chain}
summarizes the evidence chain and separates the schematic story from empirical
panels used later in the paper.

\begin{figure*}[t]
\centering
\includegraphics[width=\textwidth]{figures/fig_acl_method_chain_imagegen_v18.png}
\caption{Evidence-chain overview. We start from
USD scenes with MDL materials, convert them to composition-preserving noMDL
counterparts, compare matched renders for VLM grounding, and finally check
selected official InternNav navigation rollouts plus paired metrics. This
panel is a schematic roadmap; empirical render, material-effect, and navigation
figures later in the paper use recorded project artifacts.}
\label{fig:acl_method_chain}
\end{figure*}

The evidence is organized as four gates. First, four single-object Isaac Sim
assets screen visual and feature preservation. Second, GRScenes grounding uses
read-only source scenes, scratch no-MDL conversion, matched renders, target
projection, blind visual QA, and local Gemma4/Qwen2.5-VL probes: the 15-pair
clean pool is a pilot preservation control, while the frozen 30-pair stress pool
tests visible targets under material, color, or lighting shifts. Third, a
selected material-effect audit compares original MDL, ConvertAsset noMDL, and
the NVIDIA PreviewSurface baseline only inside bins with bounded evidence.
Fourth, a 99-episode official KuJiaLe InternNav run checks whether converted
scenes remain usable in a real navigation stack. No gate substitutes for broad
downstream robustness: not a proxy score, selected visual panel, or navigation
sanity run.


We make four contributions:
\begin{enumerate}
    \item We frame material conversion as a controlled VLM-grounding
    perturbation, not neutral asset optimization.
    \item We define a claim-bounded grounding protocol with immutable inputs,
    scratch conversion, matched renders, projection QA, prompt contracts, and
    coordinate-frame scoring.
    \item We instantiate the protocol on GRScenes with a clean preservation
    pilot and a frozen 30-pair target-centered material-shift stress set,
    showing that answer accuracy, point grounding, and coordinate compliance
    can diverge across model backends.
    \item We attach material-effect and embodied-data sanity evidence to the
    same claim registry: a selected NVIDIA-baseline material audit, a
    99-episode official KuJiaLe InternNav sanity run, and repeated official
    scene load/render stability checks. These support bounded reliability
    claims, not broad embodied-benchmark or speedup claims.
\end{enumerate}
